\chapter{Experiments and Evaluation}
\label{chap:evaluation}
\textit{This chapter evaluates \sysname{}: the research questions, the datasets and baselines, the metrics and the experimental setup, followed by the results and their discussion.}

\section{Research Questions}
\label{sec:eval-rq}
\todo{Restate RQ1--RQ3 from Section~\ref{sec:problem-statement}.}

\section{Experimental Setup}
\label{sec:eval-setup}

\subsection{Datasets}
\label{sec:eval-datasets}
\todo{Benchmark software systems, e.g.\ the cases of MARE (ATM, Cafeteria Ordering, Library, \ldots)~\cite{mare2024}, iReDev (Bookkeeping, Roll Call, \ldots)~\cite{jin2025iredev}, VibeIntent~\cite{vibeintent2026}; reference requirements; dev and test split.}

\subsection{Baselines and Variants}
\label{sec:eval-baselines}
\todo{Single-agent LLM baseline; re-implementation or reported numbers of MARE and iReDev; ablations of \sysname{}: ungrounded personas (workflow version 1) vs.\ grounded personas (version 2), without rework loop, with role-only stakeholders.}

\subsection{Evaluation Harness}
\label{sec:eval-harness}
Experiments run the same agent workflow as the product through an evaluation harness, without the web app or the database (Figure~\ref{fig:eval-harness}). A review policy answers the checkpoints (scripted answers, always approve, or an LLM reviewer), and one record is written per run with its steps, usage and artifacts.

\begin{figure}[H]
    \centering
    \resizebox{\linewidth}{!}{% UML component diagram of the evaluation harness (Section 7.2.3).
\begin{tikzpicture}[uml]
    \umlactor{res}{(-2.8,0)}{Researcher}
    \node[component] (runner) at (1.6,0) {\stereo{component}\\Runner};
    \node[component] (policy) at (1.6,2.2) {\stereo{component}\\Review policy};
    \node[component] (wf) at (5.9,0) {\stereo{component}\\Agent workflow\\\tiny imported from the backend};
    \node[component] (astore) at (4.1,-2.3) {\stereo{component}\\Artifact store\\\tiny in memory};
    \node[component] (cstore) at (7.7,-2.3) {\stereo{component}\\Checkpoint store\\\tiny in memory};
    \node[component] (metrics) at (1.6,-4.4) {\stereo{component}\\Metrics and judges};
    \begin{scope}[on background layer]
        \node[draw, thick, fill=gray!6, fit=(runner)(policy)(wf)(astore)(cstore)(metrics), inner sep=0.45cm, label={[anchor=north west, font=\scriptsize]north west:\stereo{subsystem} Harness process}] (harness) {};
    \end{scope}
    \node[artifact, text width=1.6cm] (cfg) at (-2.8,-1.9) {\stereo{artifact}\\config.yaml\\dataset.jsonl};
    \node[artifact, text width=1.6cm] (records) at (-2.8,-4.4) {\stereo{artifact}\\runs.jsonl\\artifacts/};
    \node[artifact] (summary) at (5.9,-4.4) {\stereo{artifact}\\summary.csv};
    \node[component] (llm) at (11.0,0.8) {\stereo{external system}\\LLM providers};
    \node[component] (trace) at (11.0,-1.2) {\stereo{system}\\LLM tracing tool};

    \draw[dep] (res) -- node[guard, above] {starts} (runner);
    \draw[dep] (runner.west) -- node[guard, sloped, above] {reads} (cfg.east);
    \draw[dep] (runner) -- node[guard, right] {checkpoint, decision} (policy);
    \draw[dep] (runner) -- node[guard, above, yshift=2pt] {runs each case} (wf);
    \draw[dep] (wf) -- node[guard, left] {writes} (astore);
    \draw[dep] (wf) -- node[guard, right] {saves state} (cstore);
    \draw[dep] (runner.south west) -- node[guard, sloped, above, pos=0.55] {one record per run} (records.north east);
    \draw[dep] (metrics) -- node[guard, above] {reads} (records);
    \draw[dep] (metrics) -- node[guard, above] {writes} (summary);
    \draw[dep] (wf.east) -- node[guard, above, sloped] {chat model calls} (llm.west);
    \draw[dep] (wf.east) -- node[guard, below, sloped] {traces, optional} (trace.west);
    \draw[dep] (metrics.south) -- ++(0,-0.95) -- ($(llm.east |- metrics.south)+(0.6,-0.95)$) node[guard, pos=0.5, above] {judge calls} |- (llm.east);
\end{tikzpicture}
}
    \caption{UML component diagram of the evaluation harness.}
    \label{fig:eval-harness}
\end{figure}

\subsection{Metrics}
\label{sec:eval-metrics}
\todo{Metrics of Section~\ref{sec:bg-metrics} for each research question: quality, coverage and diversity, cost, latency, reliability.}

\subsection{Implementation Details}
\label{sec:eval-details}
\todo{Models per role, temperature, run parameters, number of repeats, judge model.}

\section{Results}
\label{sec:eval-results}

\subsection{Requirement Coverage and Diversity}
\label{sec:eval-coverage}
\todo{Table: CHV, MDC, Vendi Score, number of requirements, recall against the reference, per system and variant.}

\subsection{Specification Quality}
\label{sec:eval-quality}
\todo{Table: ISO/IEC/IEEE 29148 and INVEST scores by LLM-as-a-judge and by human raters.}

\subsection{Ablation Study}
\label{sec:eval-ablation}
\todo{Effect of grounding, of the rework loop and of the number of personas.}

\subsection{Cost and Latency}
\label{sec:eval-cost}
\todo{Tokens, cost and duration per run and per agent role.}

\subsection{Case Study}
\label{sec:eval-case-study}
\todo{One run end to end: idea, personas with their evidence, an excerpt of an interview, a latent need found, the requirement and the user story it leads to, traced back.}

\section{Discussion}
\label{sec:eval-discussion}
\todo{Answer each research question; threats to validity (LLM-as-a-judge bias, benchmark leakage, non-determinism).}
